\documentclass[11pt]{article}
\usepackage[margin=1in]{geometry}
\usepackage{amsmath,amssymb,amsthm}
\usepackage{mathtools}
\usepackage[colorlinks=true,linkcolor=blue,citecolor=blue]{hyperref}
\usepackage[capitalise,nameinlink]{cleveref}
\crefname{hypothesis}{Hypothesis}{Hypotheses}
\crefname{condition}{Condition}{Conditions}
\crefname{openproblem}{Open Problem}{Open Problems}

\newtheorem{theorem}{Theorem}[section]
\newtheorem{proposition}[theorem]{Proposition}
\newtheorem{lemma}[theorem]{Lemma}
\newtheorem{corollary}[theorem]{Corollary}
\theoremstyle{definition}
\newtheorem{hypothesis}[theorem]{Hypothesis}
\newtheorem{remark}[theorem]{Remark}
\newtheorem{condition}[theorem]{Condition}
\newtheorem{openproblem}[theorem]{Open Problem}

\newcommand{\R}{\mathbb{R}}
\newcommand{\dd}{\,\mathrm{d}}
\newcommand{\Div}{\operatorname{div}}
\newcommand{\grad}{\nabla}
\emergencystretch=1.5em

\title{A bounded-class Liouville lemma for a two-dimensional steady\\
advection--diffusion equation, and the classes where it holds}
\author{Fan Li\\[8pt]
{\small School of Petroleum Engineering, China University of Petroleum (East China),}\\[2pt]
{\small Qingdao 266580, China}}
\date{}

\begin{document}
\maketitle

\begin{abstract}
We study the two-dimensional steady advection--diffusion equation
$V\cdot\nabla G=\nu\Delta_\mu G$ on the half-plane $\{\rho>0\}\times\R$, where
$\Delta_\mu=\partial_\rho^2-\rho^{-1}\partial_\rho+\partial_\zeta^2$ and $V$ is
a bounded weighted-divergence-free field,
$\partial_\rho(\rho V^\rho)+\partial_\zeta(\rho V^\zeta)=0$, the meridional
velocity of an incompressible axisymmetric field. For axis-regular
$G\in C^2\cap L^\infty$ with $G=O(\rho^2)$ we prove $G\equiv0$ in the decay
class, for shear velocity fields, for zero-advection solutions, and for the
product ansatz $G=\varphi(\rho)h(\zeta)$ over a large range of profiles: for
every positive nonconstant $2\pi$-periodic $h\in C^2(\R)$, flat minima
included and with no restriction on the order or the number of the minima,
whenever $\varphi$ satisfies two one-dimensional tail conditions. The proof
reduces the level curves of $G$ to the circle through an exact identity. The
obstruction is a sign-indefinite bulk flux with no gradient factor, and we
show that weighted incompressibility is essential by an explicit bounded
counterexample. The general bounded class remains open.
\end{abstract}

\section{Introduction}

\subsection{The equation}
Let $\nu>0$. On the half-plane $\{\rho>0\}\times\R$ consider the steady
equation
\begin{equation}\label{eq:2d-intro}
V\cdot\nabla G=\nu\Delta_\mu G,\qquad
\Delta_\mu:=\partial_\rho^2-\frac1\rho\partial_\rho+\partial_\zeta^2 ,
\end{equation}
for a bounded field $V=(V^\rho,V^\zeta)$ that is \emph{weighted
divergence-free},
\begin{equation}\label{eq:wdiv-intro}
\partial_\rho(\rho V^\rho)+\partial_\zeta(\rho V^\zeta)=0 ;
\end{equation}
this is exactly the class of meridional velocities of bounded incompressible
axisymmetric fields on $\R^3$, written in cylindrical coordinates with the
axis $\rho=0$ included. The operator $\Delta_\mu$ is the axisymmetric
Laplacian acting on scalars that are even in $\rho$; it is self-adjoint in
$L^2(\rho^{-1}d\rho\,d\zeta)$.

The equation is natural in its own right --- it is the steady advection
equation for a scalar carried by an incompressible meridional flow, at the
scaling-critical regularity of the classical Liouville theory --- and it is
also exactly the problem that the endpoint of the axisymmetric
Navier--Stokes Liouville theorem reduces to, see
\cite{companion}; the reader who wants only the two-dimensional statement may
ignore that origin entirely.

We ask: does \eqref{eq:2d-intro} admit a nonzero bounded solution in the
axisymmetric class? The answer is yes for the linear theory without
\eqref{eq:wdiv-intro} (\S\ref{sec:obstruction}) and no in each of the classes
below.

\subsection{Results}
Throughout, $G$ is assumed \emph{axis regular}: $G$ extends to a $C^2$
function of $(\rho^2,\zeta)$, equivalently $G=O(\rho^2)$ and
$\partial_\rho G=O(\rho)$ as $\rho\to0$, so that $G(0,\zeta)=0$. We write
$G=\rho^2h$ with $h$ continuous across the axis.

\begin{enumerate}
\item \textbf{Decay class} (\cref{prop:2d}). If $G\to0$ at spatial infinity
then $G\equiv0$. The proof is a Feynman--Kac/optional-stopping argument for
the generator $\nu\Delta_\mu-V\cdot\nabla$; the weight $\rho$ in $\Delta_\mu$
makes the axis an entrance boundary.
\item \textbf{Shear and zero advection} (\cref{prop:shear,prop:zeroadv}). If
$V=(0,v(\rho))$ with $v$ bounded measurable, or if $V\cdot\nabla G\equiv0$,
then $G\equiv0$; no periodicity, decay or mean condition on $v$ is used. In
the zero-advection case $G=\rho^2u$ with $u$ harmonic on $\R^5$, and
boundedness kills $u$ by a mean-value estimate.
\item \textbf{Products} (\cref{prop:product,prop:min,prop:periodic,prop:periodicgen,prop:periodicflat}).
For $G=\varphi(\rho)h(\zeta)$ the equation integrates along the level curves
of $G$, which are graphs $\rho=\rho_c(\zeta)$, and reduces to the statement
that a certain period integral $P(c)$ diverges as $c$ approaches the minimum
of $h$. We verify this for the entire class of positive nonconstant
$2\pi$-periodic $h\in C^2(\R)$ when $\varphi$ satisfies two one-dimensional
tail conditions, with no hypothesis on the minima of $h$ at all.
\end{enumerate}

\subsection{The obstruction}
The exact energy identity obtained by testing \eqref{eq:2d-intro} against
$G/\rho$ contains a bulk term $-\int V^\rho G^2/\rho^2$ that is
sign-indefinite, carries no gradient factor, and is not controlled by
$\|V\|_\infty$; the surface terms are equally uncontrolled in the bounded
class. An explicit field shows that boundedness, axis regularity and
$G=O(\rho^2)$ alone do not suffice once \eqref{eq:wdiv-intro} is dropped
(\cref{rem:2dopen}), and that the weighted incompressibility is what forces
the relevant relations to hold automatically rather than by hypothesis.

\section{The class and the notation}
Let $V=(V^\rho,V^\zeta)$ be bounded on $\{\rho>0\}\times\R$ and satisfy
\eqref{eq:wdiv-intro}; we write $\langle f\rangle_\zeta$ for a $\zeta$-mean
wherever one is taken. A function $G\in C^2\cap L^\infty$ is axis regular if
$G$ extends to a $C^2$ function of $(\rho^2,\zeta)$.

\begin{equation}
V\cdot\nabla G=\nu\Delta_\mu G\quad\text{on }\{\rho>0\}\times\R ,
\end{equation}
which is \eqref{eq:2d} below.

\paragraph{The stream function form.}
Every bounded smooth weighted-divergence-free $V$ is $V=\rho^{-1}\nabla^\perp\psi$
with $\psi=-\rho^2W/2$ and $|W|\le\|V^\zeta\|_\infty$; the equation becomes
\begin{equation}\label{eq:2dW-intro}
\frac{\rho}{2}\bigl(W_\zeta G_\rho-W_\rho G_\zeta\bigr)-WG_\zeta
=\nu\Delta_\mu G ,
\end{equation}
whose characteristic form is used in \S\ref{sec:products}.

\section{The decay class}
\begin{proposition}[2D lemma: the decay class]\label{prop:2d}
Let $G\in C^2\cap L^\infty$ satisfy
\begin{equation}\label{eq:2d}
V\cdot\nabla G=\nu\Delta_\mu G\quad\text{on }\{\rho>0\}\times\R,
\end{equation}
suppose $G$ is axis regular, $G=O(\rho^2)$ as $\rho\to0$ (equivalently, $G$
extends to a $C^2$ function of $(\rho^2,\zeta)$, i.e.\
$G(0,\zeta)=\partial_\rho G(0,\zeta)=0$), and suppose that
$G(x)\to0$ as $|(\rho,\zeta)|\to\infty$ in the half-plane
$\{\rho>0\}$. Then $G\equiv0$.
\end{proposition}
\begin{proof}
Let $L=\nu\Delta_\mu-V\cdot\nabla$ be the generator and $X$ a solution of the
associated martingale problem started at $x$ (well posed because $V$ is bounded
and the diffusion is non-degenerate). Let $\tau_S$ be the exit time of
$B_S=\{\rho^2+\zeta^2<S^2\}$ from the half-plane $\{\rho>0\}$. Let
$M=\|V^\zeta\|_\infty$, $\lambda=1+M/\nu$, and $f=e^{\lambda\zeta}$. Then
$f\in C^2$ and
\begin{equation}
Lf=\nu\lambda^2f-\lambda V^\zeta f
=\lambda f(\nu\lambda-V^\zeta)
\ge\lambda\nu e^{-\lambda S}=:c_S>0
\qquad\text{on }B_S ,
\end{equation}
because $\nu\lambda-V^\zeta\ge\nu$. Since $f$ is bounded on $B_S$, Dynkin's
formula for the stopped process gives
\begin{equation}
\mathbb E_x\bigl[f(X_{t\wedge\tau_S})\bigr]-f(x)
=\mathbb E_x\int_0^{t\wedge\tau_S}Lf(X_s)\,ds
\ge c_S\,\mathbb E_x[t\wedge\tau_S].
\end{equation}
As $0\le f\le e^{\lambda S}$ on $B_S$,
$\mathbb E_x[\tau_S]\le e^{\lambda S}/c_S<\infty$, hence $\tau_S<\infty$
a.s. Since $LG=0$ and $G$ is bounded, $G(X_{t\wedge\tau_S})$ is
a bounded martingale, and optional stopping at $\tau_S$ gives
\begin{equation}
G(x)=\mathbb E_x\bigl[G(X_{\tau_S})\bigr].
\end{equation}
At the exit point $X_{\tau_S}$ either $\rho=0$, where $G=0$ by axis regularity,
or $|(\rho,\zeta)|=S$, where $|G|\le\varepsilon(S)$ with $\varepsilon(S)\to0$
by hypothesis. Hence $|G(x)|\le\varepsilon(S)$ for every $S$, and letting
$S\to\infty$ gives $G(x)=0$.
\end{proof}
\section{Shear fields and zero advection}

\begin{proposition}[the shear class is rigid]\label{prop:shear}
Let $V=(0,v(\rho))$ with $v$ bounded and measurable, and let
$G\in C^2\cap L^\infty$ be axis regular with $G=O(\rho^2)$. If
$V\cdot\nabla G=\nu\Delta_\mu G$ on $\{\rho>0\}\times\R$, then $G\equiv0$.
No periodicity, decay at infinity, or zero mean of $v$ is used.
\end{proposition}
\begin{proof}
Let $L=\nu\Delta_\mu-V\cdot\nabla$, so $LG=0$. Since $V^\rho=0$, the
$\rho$-coordinate of the diffusion associated with $L$ is autonomous:
\begin{equation}
d\rho=-\frac{\nu}{\rho}\,dt+\sqrt{2\nu}\,dB_\rho ,
\qquad
d\zeta=-v(\rho)\,dt+\sqrt{2\nu}\,dB_\zeta .
\end{equation}
The martingale problem is well posed up to the exit time $\tau_0$ of the
axis (the coefficients are locally bounded on $\{\rho>0\}$, the $\rho$-drift
is locally Lipschitz there, the $\zeta$-drift is bounded measurable, and the
diffusion is non-degenerate). The scale density of the $\rho$-marginal is
\begin{equation}
s'(\rho)=\exp\Bigl(-\int^\rho\frac{2(-\nu/r)}{2\nu}\,dr\Bigr)=\rho ,
\end{equation}
so $s(\rho)=\rho^2/2$. For $0<\rho_0<n$,
\begin{equation}
\mathbb P_{\rho_0}(\tau_0<\tau_n)
=\frac{s(n)-s(\rho_0)}{s(n)-s(0)}
=1-\frac{\rho_0^2}{n^2}\longrightarrow1
\qquad(n\to\infty),
\end{equation}
hence $\tau_0<\infty$ a.s. Since $G$ is bounded, It\^o's formula and
$LG=0$ make $G(X_{t\wedge\tau_0})$ a bounded martingale; optional stopping
gives $G(x)=\mathbb E_x[G(X_{\tau_0})]$. At $\tau_0$ one has $\rho=0$ and
$G(0,\zeta)=0$ by axis regularity, so $G(x)=0$ for every $x$.
\end{proof}

\begin{proposition}[the zero-advection class is rigid]\label{prop:zeroadv}
Let $G\in C^2\cap L^\infty$ be axis regular with $G=O(\rho^2)$ and satisfy
\eqref{eq:2d}. Suppose $V\cdot\nabla G\equiv0$ on $\{\rho>0\}\times\R$ (the
advection hypothesis holds in particular when $G=F(\psi)$ and
$V=\rho^{-1}\nabla^\perp\psi$, with $F(0)=0$). Then $G\equiv0$.
\end{proposition}
\begin{proof}
The equation gives $\Delta_\mu G=0$. Write $G=\rho^2u$. The identity
$\Delta_\mu(\rho^2u)=\rho^2\Delta_5u$ gives $\Delta_5u=0$ on
$\{\rho>0\}$. Axis regularity makes $u=G/\rho^2$ continuous, hence locally
bounded, across the axis; since a line in $\R^5$ is removable for locally
bounded harmonic functions, $u$ is harmonic on all of $\R^5$ (and smooth
there by Weyl's lemma).
Since $|G|\le C$, we have $|u|\le C\rho^{-2}$. For $P\in\R^5$ and
$R\ge2|P|$, the mean-value property and the integrability of $\rho^{-2}$
near the axis (the axis is a line in $\R^5$, of codimension four, so that
$\int_{B_R(P)}\rho^{-2}\,dx\le CR^3$) give
\begin{equation}
|u(P)|\le\frac{C}{|B_R|}\int_{B_R(P)}\frac{dx}{\rho^2}
\le\frac{C'}{R^2}\longrightarrow0
\qquad(R\to\infty).
\end{equation}
Hence $u\equiv0$ and $G\equiv0$.
\end{proof}
\section{Products: the level-set reduction}\label{sec:products}

\begin{proposition}[the product class is rigid]\label{prop:product}
Let $\varphi\in C^2([0,\infty))$ satisfy $\varphi(0)=\varphi'(0)=0$,
$\varphi'>0$ on $(0,\infty)$, $0<\varphi<1$, and $\varphi(\rho)\to1$ as
$\rho\to\infty$. Let $h\in C^\infty(\R)$ be bounded and strictly increasing,
with $h(\zeta)\to H>0$ and $h'(\zeta)\to0$ as $\zeta\to\infty$, and
$h(-\infty)<H$. Then there is no bounded $W\in C^1$ solving
\begin{equation}\label{eq:2dW}
\frac{\rho}{2}\bigl(W_\zeta G_\rho-W_\rho G_\zeta\bigr)-WG_\zeta
=\nu\Delta_\mu G
\end{equation}
with $G=\varphi(\rho)h(\zeta)$. Consequently no bounded smooth weighted-div-free
$V$ in this class solves $V\cdot\nabla G=\nu\Delta_\mu G$ with $G\not\equiv0$.
\end{proposition}
\begin{proof}
Since $\varphi'>0$ and $h'\ne0$, $\nabla G=(\varphi'h,\varphi h')\ne0$ and the
level sets of $G$ are non-closed arcs. Fix $c\in I:=(\max\{0,h(-\infty)\},H)$,
so $c/H\in(0,1)$ and $\rho_*:=\varphi^{-1}(c/H)\in(0,\infty)$. The arc
$\gamma_c=\{G=c\}$ is parametrized by $\zeta>\zeta_c$ (where $h(\zeta_c)=c$) as
$\rho=\rho_c(\zeta):=\varphi^{-1}(c/h(\zeta))$, and
$\rho_c(\zeta)\to\rho_*$ as $\zeta\to\infty$. Set $F=\rho^2W$. Along $\gamma_c$,
\eqref{eq:2dW} gives
\begin{equation}
F'(\zeta)=2\nu\rho_c(\zeta)\Bigl(A(\rho_c(\zeta))
+k(\rho_c(\zeta))\frac{h''(\zeta)}{h(\zeta)}\Bigr),\qquad
A=\frac{\varphi''-\varphi'/\rho}{\varphi'},\quad k=\frac{\varphi}{\varphi'}>0.
\end{equation}
Let $g=h'/h$. Since $h$ is bounded, increasing, $h\to H>0$ and $h'\to0$, we
have $g\to0$, $g\in L^1$ near $+\infty$, and $g^2\le g$ eventually, so
$g^2\in L^1$. With $B_c=2\nu\rho_c k(\rho_c)\to b_*:=2\nu\rho_*k(\rho_*)>0$ and
$B_c'=O(g)$, integration by parts gives
\begin{equation}
\int_{\zeta_0}^{\zeta}B_c\frac{h''}{h}
=\int_{\zeta_0}^{\zeta}B_c(g'+g^2)
=[B_cg]_{\zeta_0}^{\zeta}-\int_{\zeta_0}^{\zeta}B_c'g
+\int_{\zeta_0}^{\zeta}B_cg^2=O(1).
\end{equation}
Since $\rho_c\to\rho_*$ and $A$ is continuous at $\rho_*>0$,
$2\nu\rho_cA(\rho_c)\to\alpha(c):=2\nu\rho_*A(\rho_*)$; hence, by Ces\`aro
averaging,
\begin{equation}
F(\zeta)=\alpha(c)\zeta+o(\zeta)+O(1).
\end{equation}
If $\alpha(c)=0$ for every $c\in I$, then $A\equiv0$ on the interval
$\rho_*(I)=(\rho_{\min},\infty)$ with $\rho_{\min}\ge0$, so
$\varphi''=\varphi'/\rho$ and $\varphi'=K\rho$ there; then
$\varphi=K\rho^2/2+C$ on that interval, and boundedness with $0<\varphi<1$ and
$\varphi\to1$ forces $K=0$ and $\varphi\equiv1$ there, a contradiction. Hence
there is $c\in I$ with $\alpha(c)\ne0$. For this $c$,
$F(\zeta)=\alpha(c)\zeta+o(\zeta)\to\pm\infty$. Since
$\rho_c(\zeta)\to\rho_*\in(0,\infty)$, we get $|W|=|F|/\rho_c^2\to\infty$,
contradicting the boundedness of $W$.

Finally, if $V$ is bounded, smooth and weighted divergence-free, then
$\psi=\int_0^\rho rV^\zeta(r,\zeta)\,dr$ and $W=-2\psi/\rho^2$ give a bounded
$C^1$ solution of \eqref{eq:2dW}; so no such $V$ solves
$V\cdot\nabla G=\nu\Delta_\mu G$ with $G$ in this class.
\end{proof}

\begin{proposition}[product class with a nondegenerate minimum]\label{prop:min}
Let $\nu>0$ and let $\varphi\in C^2([0,\infty))$ satisfy
$\varphi(0)=\varphi'(0)=0$,
$\varphi'>0$ on $(0,\infty)$, $0<\varphi<1$, and $\varphi(\rho)\to1$ as
$\rho\to\infty$. Let $h\in C^2(\R)$ be positive with a nondegenerate minimum at
$\zeta_0$: $h'(\zeta_0)=0$, $\beta:=h''(\zeta_0)>0$, $m:=h(\zeta_0)>0$. Then no
bounded smooth weighted-div-free $V$ solves $V\cdot\nabla G=\nu\Delta_\mu G$
with $G=\varphi(\rho)h(\zeta)$ and $G\not\equiv0$; in fact, there is no
$W\in C^1$ with $W$, $\rho W_\rho$, $\rho W_\zeta$ bounded solving
\eqref{eq:2dW}.
\end{proposition}
\begin{proof}
Let $W\in C^1$ with $W,\rho W_\rho,\rho W_\zeta$ bounded solve
\eqref{eq:2dW}; a bounded smooth weighted-div-free $V$ yields such a $W$ by the
representation at the end of the proof of \cref{prop:product}. Put
$V^\rho=\rho W_\zeta/2$.
At $\zeta=\zeta_0$ we have $G_\zeta=\varphi h'=0$ and $G_\rho=\varphi'h>0$.
Equation \eqref{eq:2dW} gives
\begin{equation}
V^\rho=\frac{\rho}{2}W_\zeta
=\nu\frac{\Delta_\mu G}{G_\rho}
=\nu\Bigl(A(\rho)+k(\rho)\frac{h''(\zeta_0)}{h(\zeta_0)}\Bigr),
\quad A=\frac{\varphi''-\varphi'/\rho}{\varphi'},\quad
k=\frac{\varphi}{\varphi'}>0.
\end{equation}
Set $\omega^2=\beta/m>0$. Since $V^\rho$ is bounded, there is $C$ with
\begin{equation}
\Bigl|\frac{\varphi''-\varphi'/\rho}{\varphi'}
+\omega^2\frac{\varphi}{\varphi'}\Bigr|\le C,
\end{equation}
hence $|\varphi''-\varphi'/\rho+\omega^2\varphi|\le C\varphi'$ and therefore
\begin{equation}
\varphi''+\omega^2\varphi\le\Bigl(C+\frac1\rho\Bigr)\varphi'.
\end{equation}
Integrating from $\rho_0>0$ to $R$,
\begin{equation}
\varphi'(R)\le\varphi'(\rho_0)+C\bigl(\varphi(R)-\varphi(\rho_0)\bigr)
+\int_{\rho_0}^R\frac{\varphi'}{\rho}\,d\rho
-\omega^2\int_{\rho_0}^R\varphi\,d\rho .
\end{equation}
The first three terms on the right are bounded: $\varphi$ is bounded,
$\varphi'\in L^1(0,\infty)$ because $\varphi$ is bounded and increasing, and
$\varphi'/\rho$ is integrable near $0$ (as $\varphi'(0)=0$ and
$\varphi\in C^2$) and bounded by $\varphi'$ for $\rho\ge1$. Since
$\varphi\to1$, the last term tends to $-\infty$, so $\varphi'(R)<0$ for large
$R$, contradicting $\varphi'>0$. This rules out every such $W$ and hence, by
the representation, every bounded smooth weighted-div-free $V$. The sign
$\omega^2=\beta/m>0$ is
essential: at a nondegenerate maximum the decisive term in the integrated
inequality flips sign, and this argument gives no contradiction.
\end{proof}

\begin{remark}[periodic level graphs and the period integral]\label{rem:drift}
For $2\pi$-periodic $h$ and $c<\min h$ the level graph
$\rho_c(\zeta)=\varphi^{-1}(c/h(\zeta))$ is
periodic in $\zeta$ but is not a closed curve in the plane, so the closed-loop
identity \eqref{eq:2dfredholm} does not apply; the period integral $P(c)$
defined below need not vanish. Set
$F(\zeta)=\rho_c(\zeta)^2W(\rho_c(\zeta),\zeta)$. For
$\varphi=\rho^2/(1+\rho^2)$ and $h=1+\varepsilon\cos\zeta$, $0<\varepsilon<1$,
a direct computation gives
\begin{equation}
F'=\nu c\Bigl(\frac{h''}{(h-c)^2}-\frac{8}{h}\Bigr),\qquad
P(c):=\int_0^{2\pi}F'\,d\zeta
=2\pi\nu c\Bigl(\frac{\varepsilon^2}{((1-c)^2-\varepsilon^2)^{3/2}}
-\frac{8}{\sqrt{1-\varepsilon^2}}\Bigr).
\end{equation}
Thus $P(c)$ is not identically zero (the bracket is nonconstant in $c$, and
$P(c)\to+\infty$ as $c\uparrow1-\varepsilon$). If a bounded $W$ existed, then
$F=\rho_c^2W$ would be bounded on the periodic level graph and
$F(\zeta+2\pi n)=F(\zeta)+nP(c)$ for every $n$, so $P(c)=0$ for every
$c<\min h$, a contradiction. Explicitly, for every $c$ with $P(c)\neq0$ the formal solution
$W=[\int_0^\zeta R(G,s)\,ds+H(G)]/\rho^2$, with
$R(c,s)=2\nu\rho_c(s)[A(\rho_c(s))+k(\rho_c(s))h''(s)/h(s)]$, grows linearly in
$\zeta$ along that level graph. In particular, for $h=1+\varepsilon\cos\zeta$
the minimum at $\zeta_0=\pi$ has $m=1-\varepsilon>0$ and $\beta=\varepsilon>0$,
so \cref{prop:min} applies for every $\varphi$ satisfying its hypotheses. The
period integral gives a second, explicit proof in this example. Periodicity is
assumed in this remark and in \cref{prop:periodic}; \cref{prop:min} does not
use it.
\end{remark}

\section{Periodic products}
\begin{proposition}[the periodic product class for $\varphi_1$]\label{prop:periodic}
Let $\nu>0$, $\varphi=\varphi_1=\rho^2/(1+\rho^2)$, and let $h\in C^2(\R)$ be
positive and $2\pi$-periodic. Then there is no bounded $W\in C^1$ solving
\eqref{eq:2dW} with $G=\varphi_1(\rho)h(\zeta)$; consequently no bounded smooth
weighted-div-free $V$ solves $V\cdot\nabla G=\nu\Delta_\mu G$ with this $G$ and
$G\not\equiv0$.
\end{proposition}
\begin{proof}
First suppose $h$ is constant. For $c\in(0,h)$ the level graph
$\rho_c\equiv\sqrt{c/(h-c)}$ is a circle. Since $G_\zeta=0$, \eqref{eq:2dW}
gives $(\rho/2)W_\zeta G_\rho=\nu\Delta_\mu G$, hence
$W_\zeta=2\nu\Delta_\mu G/(\rho G_\rho)=-8\nu/(1+\rho^2)$ and
$F'=\rho_c^2W_\zeta=-8\nu c/h$, a nonzero constant; then
$F(\zeta+2\pi n)=F(\zeta)-16\pi\nu c n/h$ for every $n$, so a bounded $W$ is
impossible.

Now suppose $h$ is nonconstant, and put $m=\min h$, $M=\max h$, so
$0<m<M$. Fix $c\in(0,m)$, set $\delta=m-c$ and $u=h-c>0$. The level graph
$\rho_c(\zeta)=\sqrt{c/u(\zeta)}$ is $2\pi$-periodic and satisfies
$\varphi_1(\rho_c)h=c$. Set $F(\zeta)=\rho_c(\zeta)^2W(\rho_c(\zeta),\zeta)$.
Along this graph, \eqref{eq:2dW} gives
\begin{equation}
F'=2\nu\rho_c\Bigl(A(\rho_c)+k(\rho_c)\frac{h''}{h}\Bigr),
\qquad A=\frac{\varphi_1''-\varphi_1'/\rho}{\varphi_1'},\quad
k=\frac{\varphi_1}{\varphi_1'}>0.
\end{equation}
For $\varphi_1$ one has $A=-4\rho/(1+\rho^2)$ and
$k=\rho(1+\rho^2)/2$; since $\rho_c^2=c/u$ and
$1+\rho_c^2=h/u$, this becomes
\begin{equation}\label{eq:periodicF}
F'=\nu c\Bigl(\frac{h''}{(h-c)^2}-\frac8h\Bigr)=:f(\zeta),
\end{equation}
which is $2\pi$-periodic. Let $P(c)=\int_0^{2\pi}f\,d\zeta$. Since $u=h-c$ is
positive, $2\pi$-periodic and $C^2$, integration by parts on the circle gives
\begin{equation}\label{eq:periodicI}
I(c):=\int_0^{2\pi}\frac{h''}{(h-c)^2}\,d\zeta
=2\int_0^{2\pi}\frac{(h')^2}{(h-c)^3}\,d\zeta\ge0,
\end{equation}
so
\begin{equation}
P(c)=\nu c\,I(c)-8\nu c\int_0^{2\pi}\frac{d\zeta}{h(\zeta)}.
\end{equation}
The second term is bounded: $h\ge m$ gives
$8\nu c\int_0^{2\pi}d\zeta/h\le16\pi\nu$. To bound $I(c)$ from below, put
$\psi=u^{-1/2}\in C^1$. Then $\psi'=-\frac12h'u^{-3/2}$, hence
$I(c)=8\int_0^{2\pi}(\psi')^2\,d\zeta$. The function $u$ ranges over
$[\delta,M-m+\delta]$, so the oscillation of $\psi=u^{-1/2}$ is
\begin{equation}
\operatorname{osc}\psi=\delta^{-1/2}-(M-m+\delta)^{-1/2}.
\end{equation}
For $\delta\le(M-m)/3$ this is at least $\delta^{-1/2}/2$. The Poincar\'e
inequality on the circle,
$\int_0^{2\pi}(\psi')^2\ge(\operatorname{osc}\psi)^2/(2\pi)$, gives
\begin{equation}
I(c)\ge\frac{\delta^{-1}}{\pi}.
\end{equation}
Therefore
\begin{equation}
P(c)\ge\frac{\nu c}{\pi\delta}-16\pi\nu\longrightarrow+\infty
\qquad(c\uparrow m).
\end{equation}
Choose $c<m$ with $P(c)>0$. If a bounded $W$ existed, then
$F=\rho_c^2W$ would be bounded, because
$\rho_c^2=c/(h-c)\le c/\delta$, and $2\pi$-periodic; but $f$ is
$2\pi$-periodic, so $F(\zeta+2\pi n)=F(\zeta)+nP(c)$ for every $n$, which is
unbounded. This contradiction rules out every bounded $W$. The passage to $V$
is the representation at the end of the proof of \cref{prop:product}.
\end{proof}

\begin{proposition}[saturating periodic products with finite-order minima]\label{prop:periodicgen}
Let $\nu>0$, and let $\varphi\in C^2([0,\infty))$ satisfy
$\varphi(0)=\varphi'(0)=0$, $\varphi'>0$ on $(0,\infty)$, $0<\varphi<1$, and
$\varphi(\rho)\to1$ as $\rho\to\infty$. Put $\varepsilon=1-\varphi$ and
\begin{equation}
A=\frac{\varphi''-\varphi'/\rho}{\varphi'},\qquad k=\frac{\varphi}{\varphi'}>0 .
\end{equation}
Let $K\ge1$ be an integer, and let $h\in C^{2K+1}(\R)$ be positive and
$2\pi$-periodic, with finitely many global minima $m=h(\zeta_i)$,
$i=0,\dots,N$, each of exact order $2K_i$ with $1\le K_i\le K$, and at least
one of them of order $2K$:
\begin{equation}
h(\zeta_i+s)=m+a_i|s|^{2K_i}+O(|s|^{2K_i+1}),\qquad a_i>0 .
\end{equation}
Put $\theta=1/(2K)$, and assume there are $\rho_0\ge0$ and a nondecreasing
$g:[\rho_0,\infty)\to[0,\infty)$ with $(\rho A(\rho))_-\le g(\rho)$ and
\begin{equation}\label{eq:Htheta}
\int_{\rho_0}^{\infty}g(\rho)\,
\varepsilon(\rho)^{\theta-1}\varphi'(\rho)\,d\rho<\infty .
\end{equation}
Let $\rho_c(\zeta):=\varphi^{-1}(c/h(\zeta))$ and
\begin{equation}
P(c):=2\nu\int_0^{2\pi}\rho_c\Bigl(A(\rho_c)+k(\rho_c)\frac{h''}{h}\Bigr)d\zeta .
\end{equation}
Then $P(c)\to+\infty$ as $c\uparrow m$. Consequently there is no bounded
$W\in C^1$ solving \eqref{eq:2dW} with $G=\varphi(\rho)h(\zeta)$, and no
bounded smooth weighted-divergence-free $V$ solves
$V\cdot\nabla G=\nu\Delta_\mu G$ with this $G$.
\end{proposition}
\begin{proof}
The function $\varphi$ is strictly increasing from $0$ to $1$, so
$\Phi:=\varphi^{-1}$ is well defined on $(0,1)$. Put
$\varepsilon(\rho)=1-\varphi(\rho)$ and $r(\varepsilon)=\Phi(1-\varepsilon)$.
Let $m=\min h$, $M=\max h$, fix $c=m-\delta$ with
$0<\delta<\min(m/2,(M-m)/3)$, and set
\begin{equation}
\rho_c(\zeta)=\Phi\Bigl(\frac{c}{h(\zeta)}\Bigr).
\end{equation}
Then $G(\rho_c(\zeta),\zeta)=c$, and $\rho_c$ is $2\pi$-periodic. By the
computation in the proof of \cref{prop:product}, along this level graph
\begin{equation}\label{eq:Pgen}
P(c):=\int_0^{2\pi}F'\,d\zeta
=2\nu\int_0^{2\pi}\rho_c\Bigl(A(\rho_c)+k(\rho_c)\frac{h''}{h}\Bigr)d\zeta
=:I_1(c)+I_2(c),
\end{equation}
where $F(\zeta)=\rho_c(\zeta)^2W(\rho_c(\zeta),\zeta)$.

\emph{Step 1: $I_2(c)\to+\infty$.}  Choose a global minimum $\zeta_0$ of
maximal order $2K$.  Fix $\zeta_1>0$ so small that the intervals
$(\zeta_i-\zeta_1,\zeta_i+\zeta_1)$, $i=0,\dots,N$, are disjoint and
$\eta:=\min\{h(\zeta)-m:\ |\zeta-\zeta_i|\ge\zeta_1\ \text{for all }i\}>0$.
Then off the union of these intervals $h-c\ge\eta$, $\rho_c$ stays in a fixed
compact interval and the contribution is $O(1)$ uniformly in $c$.  The
remaining intervals $|\zeta-\zeta_i|<\zeta_1$, $i\ge1$, contribute a
nonnegative amount, since there $h''/h\ge c|\zeta-\zeta_i|^{2K_i-2}\ge0$.
Put
$W_{1,i}:=\max_{|s|\le\zeta_1}(h(\zeta_i+s)-m)$ and $W_1:=W_{1,0}$.  On
$|\zeta-\zeta_0|<\zeta_1$ the expansion of $h$ gives
$h''(\zeta)/h(\zeta)\ge c_0|\zeta-\zeta_0|^{2K-2}$.  Put
$w=h-c\in[\delta,W_1+\delta]$, and note that in the estimates below the
integration range may be restricted to $w\le W_1$, all integrands being of one
sign.  Put $\varepsilon=w/(c+w)$; then
$\rho_c=r(\varepsilon)$ and
\begin{equation}
\rho_c k(\rho_c)=\varphi(\rho_c)\frac{\rho_c}{\varphi'(\rho_c)}
=\varphi(\rho_c)\Psi(\varepsilon),\qquad
\Psi(\varepsilon):=\frac{r(\varepsilon)}{\varphi'(r(\varepsilon))}>0 .
\end{equation}
Since $\varphi(\rho_c)=c/h\ge c/M\ge m/(2M)>0$ for $c>m/2$, and since on each
side of $\zeta_0$ one has
$|\zeta-\zeta_0|^{2K-2}d\zeta
=(2K)^{-1}a^{\theta-1}(1+O(|\zeta-\zeta_0|))(w-\delta)^{-\theta}dw$, we get
\begin{equation}
I_2(c)\ge C\int_{\delta}^{W_1}\Psi\Bigl(\frac{w}{c+w}\Bigr)
(w-\delta)^{-\theta}\,dw-C_1
\end{equation}
with $C>0$ and $C_1\ge0$ independent of $c$.  On $w\ge2\delta$ we have
$(w-\delta)^{-\theta}\ge2^{-\theta}w^{-\theta}$; applying Fatou's lemma as
$\delta\downarrow0$,
\begin{equation}
\liminf_{c\uparrow m}I_2(c)\ge
C2^{-\theta}\int_0^{W_1}\Psi\Bigl(\frac{w}{m+w}\Bigr)w^{-\theta}\,dw .
\end{equation}
The substitution $w=mv/(1-v)$ gives $dw=m(1-v)^{-2}dv$ and
\begin{equation}
\int_0^{W_1}\Psi\Bigl(\frac{w}{m+w}\Bigr)w^{-\theta}dw
=m^{1-\theta}\int_0^{v_1}\Psi(v)v^{-\theta}(1-v)^{\theta-2}dv .
\end{equation}
Since $(1-v)^{\theta-2}\ge c_3>0$ on $[0,v_1]$, and since
$\Psi(v)\,dv=\bigl(r(v)/\varphi'(r(v))\bigr)d\varepsilon=-r(v)\,d\rho$
with $d\varepsilon=-\varphi'(\rho)\,d\rho$, we obtain
\begin{equation}
\int_0^{v_1}\Psi(v)v^{-\theta}dv
=\int_{r(v_1)}^{\infty}\rho\,\varepsilon(\rho)^{-\theta}\,d\rho=+\infty,
\end{equation}
because $\varepsilon(\rho)\to0$ while $\rho\ge r(v_1)>0$.  Hence
$I_2(c)\to+\infty$.

\emph{Step 2: $I_1$ is bounded below.}  With $\eta$ as in Step 1, for
$\delta<\eta/2$ one has $h-c\ge\eta/2$ off the union of the intervals
$(\zeta_i-\zeta_1,\zeta_i+\zeta_1)$, so the contribution to $I_1$ there is
$O(1)$.  Near a minimum $\zeta_i\ne\zeta_0$ of order $2K_i$ the same change of
variables carries the weight $(w-\delta)^{\theta_i-1}$ with
$\theta_i=1/(2K_i)\ge\theta$; since $\varepsilon^{\theta_i-1}\le
\varepsilon^{\theta-1}$ for $0<\varepsilon<1$, \eqref{eq:Htheta} bounds the
negative part of that contribution by the argument below with $\theta_i$ in
place of $\theta$.  It therefore suffices to treat $\zeta_0$.  Shrinking
$\zeta_1$ (hence all the $W_{1,i}$) if necessary, assume in addition that
$\Phi(m/(m+W_{1,i}))\ge\rho_0$ for every $i$; then for each $i$ and every
$0<w<W_{1,i}$ the point $\rho_m(w)$ below lies in $[\rho_0,\infty)$.  With
$W_1=W_{1,0}$ one has
$\rho_c=\Phi(c/(c+w))\ge\Phi(c/(c+W_1))\to\Phi(m/(m+W_1))$ and
$\rho_m(w)\ge\Phi(m/(m+W_1))$, so for $c$ sufficiently close to $m$ every
argument of $g$ below lies in $[\rho_0,\infty)$.  On
$|\zeta-\zeta_0|<\zeta_1$ the change of variables above, which for the
$A$-term carries no extra factor $h''/h$, gives
\begin{equation}
I_1(c)\ge -C\int_\delta^{W_1}
\bigl(\rho_c A(\rho_c)\bigr)_- (w-\delta)^{\theta-1}\,dw+O(1).
\end{equation}
For $c\le m$ one has
\begin{equation}
\frac{c}{c+w}\le\frac{m}{m+w},
\qquad\text{hence}\qquad
\rho_c\le \rho_m(w):=\Phi\Bigl(\frac{m}{m+w}\Bigr).
\end{equation}
Since $g$ is nondecreasing and $(\rho A)_-\le g$,
\begin{equation}
\bigl(\rho_c A(\rho_c)\bigr)_-\le g(\rho_c)\le g(\rho_m(w)).
\end{equation}
Therefore
\begin{equation}
\int_\delta^{W_1}\bigl(\rho_c A(\rho_c)\bigr)_-(w-\delta)^{\theta-1}dw
\le\int_\delta^{W_1}g(\rho_m(w))(w-\delta)^{\theta-1}dw .
\end{equation}
Split the last integral at $2\delta$.  On $(2\delta,W_1)$ one has
$(w-\delta)^{\theta-1}\le2^{1-\theta}w^{\theta-1}$, and
\eqref{eq:Htheta} together with the substitution $w=mv/(1-v)$,
$v=\varepsilon(\rho_m(w))=w/(m+w)$ shows that $g(\rho_m(w))w^{\theta-1}$ is
integrable on $(0,W_1)$: indeed $dw=m(1-v)^{-2}dv$,
$w^{\theta-1}=(mv/(1-v))^{\theta-1}$ and $(1-v)^{-1-\theta}\le C$ on
$[0,v_1]$, so the integral is at most
$C\int_0^{v_1}g(r(v))v^{\theta-1}dv
=C\int g(\rho)\varepsilon(\rho)^{\theta-1}\varphi'(\rho)\,d\rho<\infty$
by \eqref{eq:Htheta}.  Hence this part is at most the full integral
$\int_0^{W_1}g(\rho_m(w))w^{\theta-1}dw$, which is finite, and in particular
it is bounded uniformly in $\delta$.  On $(\delta,2\delta)$, with
$u=w-\delta\in(0,\delta)$
and $\varepsilon_\delta:=\delta/(m+\delta)$,
\begin{equation}
\int_\delta^{2\delta}g(\rho_m(w))(w-\delta)^{\theta-1}dw
\le g\bigl(\Phi(1-\varepsilon_\delta)\bigr)\frac{\delta^\theta}{\theta}
\le C\,\varepsilon_\delta^{\theta}g\bigl(\Phi(1-\varepsilon_\delta)\bigr),
\end{equation}
so this part tends to $0$ provided
$G(\varepsilon):=g(\Phi(1-\varepsilon))$ satisfies
$G(\varepsilon)\varepsilon^{\theta}\to0$.  Since $g$ is nondecreasing, $G$ is
nonincreasing, and \eqref{eq:Htheta} says
$\int_0G(\varepsilon)\varepsilon^{\theta-1}d\varepsilon<\infty$; hence
\begin{equation}
G(\varepsilon)\varepsilon^{\theta}
\le\frac{\theta}{1-2^{-\theta}}
\int_{\varepsilon/2}^{\varepsilon}G(s)s^{\theta-1}\,ds
\longrightarrow0
\qquad(\varepsilon\downarrow0).
\end{equation}
Therefore
\begin{equation}
\sup_{0<\delta<\delta_0}
\int_\delta^{W_1}g(\rho_m(w))(w-\delta)^{\theta-1}dw<\infty .
\end{equation}
Thus $I_1(c)$ is bounded below uniformly in $c$.

Combining the two steps with \eqref{eq:Pgen} gives $P(c)\to+\infty$ as
$c\uparrow m$.  In particular $P(c)>0$ for all $c$ sufficiently close to
$m$.  If a bounded $W\in C^1$ existed, then
$|F|\le\Phi(c/m)^2\|W\|_\infty$ would be bounded, while the $2\pi$-periodicity
of the integrand $F'$ in \eqref{eq:Pgen} gives
$F(\zeta+2\pi n)=F(\zeta)+nP(c)$ for every $n$; this is impossible when
$P(c)\neq0$.  The passage to a bounded weighted-divergence-free $V$ is the
same as at the end of the proof of \cref{prop:product}.
\end{proof}

\begin{remark}[explicit classes in \cref{prop:periodicgen}]\label{rem:periodicgen}
Condition \eqref{eq:Htheta} holds whenever $(\rho A)_-$ is bounded, since
then $g$ may be taken constant and
$\int_{\rho_0}^\infty\varepsilon^{\theta-1}\varphi'\,d\rho
=\int_0^{\varepsilon(\rho_0)}\varepsilon^{\theta-1}\,d\varepsilon<\infty$
(as $\theta>0$).
Condition \eqref{eq:Htheta} is sufficient, not necessary.
This covers $\varphi_1$ (for which $\rho A=-4\rho^2/(1+\rho^2)\to-4$) and
every admissible $\varphi$ with $1-\varphi\sim C\rho^{-p}$ as $\rho\to\infty$
(for instance $\varphi_p=1-(1+\rho^2)^{-p/2}$), for which $\rho A\to-(p+2)$.  It also covers profiles for which $(\rho A)_-$ grows at
most polynomially while $\varepsilon$ decays faster than any power, such as
$\varphi=1-e^{-\kappa\rho}(1+\kappa\rho)$ (with $A\equiv-\kappa$) and
stretched-exponential tails $\varepsilon\sim e^{-\rho^\alpha}$ with
$\alpha\ge2$ (for $0<\alpha<2$ one first modifies the profile near $\rho=0$,
which is not $C^2$ there, leaving the tail and \eqref{eq:Htheta} unchanged).
The periodic product class for an arbitrary admissible $\varphi$ remains
open; \cref{prop:periodicflat} extends the conclusion to a different, partly
overlapping class of profiles, with no hypothesis on $h$ beyond $C^2$.
\end{remark}

\begin{proposition}[saturating periodic products: $C^2$ data]\label{prop:periodicflat}
Let $\nu>0$, and let $\varphi\in C^2([0,\infty))$ satisfy
$\varphi(0)=\varphi'(0)=0$, $\varphi'>0$ on $(0,\infty)$, $0<\varphi<1$, and
$\varphi(\rho)\to1$ as $\rho\to\infty$. With $\Phi=\varphi^{-1}$ put
\begin{equation}
\varepsilon=1-\varphi,\qquad r(\varepsilon)=\Phi(1-\varepsilon),\qquad
\Psi(\varepsilon)=\frac{r(\varepsilon)}{\varphi'(r(\varepsilon))},
\end{equation}
and, with $A=(\varphi''-\varphi'/\rho)/\varphi'$,
\begin{equation}
\mathcal W(\varepsilon):=(1-\varepsilon)\Psi(\varepsilon),\qquad
\mathcal A(\varepsilon):=\bigl(\rho A(\rho)\bigr)_-\big|_{\rho=r(\varepsilon)}
\end{equation}
(so that $\mathcal W(\varepsilon)=\rho k(\rho)$ at $\rho=r(\varepsilon)$). Assume
\begin{enumerate}
\item[(i)] there are $\eta\in(0,2]$, $\varepsilon_1\in(0,1)$ and $C\ge0$ with
$\mathcal A(\varepsilon)\le C\bigl(1+\log\frac1\varepsilon\bigr)^{2-\eta}$
for $0<\varepsilon<\varepsilon_1$;
\item[(ii)] there are $\varepsilon_0\in(0,1)$, $c_*>0$ and $q\in[0,\eta)$
(with $\eta$ from (i)) with
$\mathcal W'(\varepsilon)\le-c_*\varepsilon^{-2}\bigl(\log\frac1\varepsilon\bigr)^{-q}$
for $0<\varepsilon<\varepsilon_0$ (differentiability of $\mathcal W$ being
automatic).
\end{enumerate}
Then for every positive nonconstant $2\pi$-periodic $h\in C^2(\R)$, with
$m=\min h$, $\rho_c=\Phi(c/h)$ and
\begin{equation}
P(c)=2\nu\int_0^{2\pi}\rho_c\Bigl(A(\rho_c)+k(\rho_c)\frac{h''}{h}\Bigr)d\zeta,
\end{equation}
one has $P(c)\to+\infty$ as $c\uparrow m$. Consequently there is no bounded
$W\in C^1$ solving \eqref{eq:2dW} with $G=\varphi(\rho)h(\zeta)$, and no
bounded smooth weighted-divergence-free $V$ solves
$V\cdot\nabla G=\nu\Delta_\mu G$ with this $G$.
\end{proposition}
\begin{proof}
Since $\varphi\in C^2$ and $\varphi'>0$ on $(0,\infty)$, $\Phi\in C^2((0,1))$
and $\mathcal W\in C^1((0,1))$. Let $m=\min h$, $M=\max h$, and fix
$c=m-\delta$ with $0<\delta<\min(m/2,(M-m)/2)$. Put
\begin{equation}
w=h-c\in[\delta,M-m+\delta],\qquad
\varepsilon(\zeta)=\frac{w(\zeta)}{h(\zeta)}=1-\frac ch,\qquad
\rho_c(\zeta)=r(\varepsilon(\zeta))=\Phi\Bigl(\frac{c}{h(\zeta)}\Bigr).
\end{equation}
Then $G(\rho_c(\zeta),\zeta)=c$, and by the computation in the proof of
\cref{prop:product}, $P=I_1+I_2$ with
\begin{equation}
I_1=2\nu\int_0^{2\pi}\rho_cA(\rho_c)\,d\zeta,\qquad
I_2=2\nu\int_0^{2\pi}\rho_ck(\rho_c)\frac{h''}{h}\,d\zeta .
\end{equation}
\emph{The $A$-term.}  Since $c\le h\le M$ and $h-c=w\ge\delta$, one has
$\varepsilon(\zeta)=w/h\ge\delta/M$ for every $\zeta$, hence
$\rho_c(\zeta)\le r(\delta/M)$.  By (i),
$\mathcal A(\varepsilon(\zeta))\le C(1+\log(M/\delta))^{2-\eta}$ on the set
$\{\varepsilon<\varepsilon_1\}$; on its complement $\varepsilon$ lies in the
compact subinterval $[\varepsilon_1,1-m/(2M)]$ of $(0,1)$ (here $c\ge m/2$,
since $\delta<m/2$), where $\mathcal A$ is continuous and hence bounded.
Enlarging $C$ if necessary,
\begin{equation}
\bigl(\rho_cA(\rho_c)\bigr)_-\le C\Bigl(1+\log\frac M\delta\Bigr)^{2-\eta}
\qquad\text{for every }\zeta,
\end{equation}
and therefore $I_1\ge-4\pi\nu C(1+\log(M/\delta))^{2-\eta}$.

\emph{The curvature term.}  Since $\varphi(\rho_c)=c/h=1-\varepsilon$,
\begin{equation}
\rho_ck(\rho_c)=\varphi(\rho_c)\frac{r(\varepsilon)}{\varphi'(r(\varepsilon))}
=(1-\varepsilon)\Psi(\varepsilon)=\mathcal W(\varepsilon).
\end{equation}
Put $v=\log h$, so that $h''/h=v''+(v')^2$ and $v'=h'/h$. The map
$\mathcal W(\varepsilon(\cdot))$ is $C^1$, and $v'$ is $2\pi$-periodic; hence so
is the product $\mathcal W(\varepsilon(\cdot))v'$, and integration by parts on
the circle gives
\begin{equation}
\int_0^{2\pi}\mathcal W(\varepsilon)v''\,d\zeta
=-\int_0^{2\pi}\mathcal W'(\varepsilon)\varepsilon'v'\,d\zeta,
\end{equation}
and the second contribution,
$2\nu\int_0^{2\pi}\mathcal W(\varepsilon)(v')^2d\zeta$, is nonnegative because
$\mathcal W>0$. Now $\varepsilon'=ch'/h^2$, hence
$\varepsilon'v'=c(h')^2/h^3\ge0$. Let
$E=\{\zeta:\varepsilon(\zeta)<\varepsilon_0\}$. On $E$ one has
$\varepsilon\ge\delta/M$, hence $\log\frac1\varepsilon\le\log\frac M\delta$ and
condition (ii) gives
\begin{equation}
-\mathcal W'(\varepsilon)\varepsilon'v'\ge
c_*\Bigl(\log\frac M\delta\Bigr)^{-q}
c\,\varepsilon^{-2}\frac{(h')^2}{h^3}.
\end{equation}
Since $m\le h\le M$ and $h'\equiv w'$, we have $\varepsilon=w/h\le w/m$, so
$\varepsilon^{-2}\ge m^2/w^2$, and $(h')^2/h^3\ge(w')^2/M^3$; therefore
\begin{equation}
-\mathcal W'(\varepsilon)\varepsilon'v'\ge\frac{c_*cm^2}{M^3}
\Bigl(\log\frac M\delta\Bigr)^{-q}\bigl((\log w)'\bigr)^2
\qquad\text{on }E .
\end{equation}
Off $E$ one has $\varepsilon\in[\varepsilon_0,1-c/M]$, where $\mathcal W'$ and
$\varepsilon'v'$ are bounded, so there is $C_1$ with
\begin{equation}
I_2\ge\frac{2\nu c_*cm^2}{M^3}\Bigl(\log\frac M\delta\Bigr)^{-q}
\int_E\bigl((\log w)'\bigr)^2d\zeta-C_1 .
\end{equation}
By Cauchy--Schwarz,
\begin{equation}
\int_E\bigl((\log w)'\bigr)^2d\zeta
\ge\frac1{2\pi}\Bigl(\int_E\bigl|(\log w)'\bigr|d\zeta\Bigr)^2
=\frac1{2\pi}\bigl(\operatorname{TV}_E\log w\bigr)^2 .
\end{equation}
Now $E=\{\varepsilon<\varepsilon_0\}=\{w<\kappa\}$ with
$\kappa=\varepsilon_0c/(1-\varepsilon_0)>0$, so $E$ contains a neighbourhood of
the set of global minima of $h$. Let $\zeta_0$ be a global minimum and put
$\eta_0:=\tfrac12\min\{\kappa,M-m+\delta\}>0$. Since
$\max w=M-m+\delta\ge2\eta_0$, one may move from $\zeta_0$ towards a point
where $w$ is maximal until $w$
first reaches $\eta_0$; on the way $w<\eta_0<\kappa$, so the path lies in $E$
and carries
$\log w$ from $\log\delta$ to $\log\eta_0$; hence
$\operatorname{TV}_E\log w\ge\log(\eta_0/\delta)$. Therefore, for
$\delta<\kappa/2$,
\begin{equation}
P(c)\ge\frac{\nu c_*cm^2}{\pi M^3}
\Bigl(\log\frac M\delta\Bigr)^{-q}\log^2\frac{\eta_0}{\delta}
-4\pi\nu C\Bigl(1+\log\frac M\delta\Bigr)^{2-\eta}-C_1
\longrightarrow+\infty\qquad(c\uparrow m),
\end{equation}
because $\eta_0\ge\tfrac12\min\{\kappa,M-m\}\ge
\tfrac12\min\{\varepsilon_0m/(2(1-\varepsilon_0)),M-m\}>0$ for $c\ge m/2$ is
bounded below independently of $\delta$, so that
$\log(\eta_0/\delta)\ge\tfrac12\log(M/\delta)$ for small $\delta$: with
$\Lambda=\log(M/\delta)$ the first term is at least
$(\nu c_*cm^2/4\pi M^3)\Lambda^{2-q}$, the second is $O(\Lambda^{2-\eta})$,
and $q<\eta$.

Choose then $c<m$ with $P(c)>0$. For this fixed $c$ one has
$\rho_c\le\Phi(c/m)<\infty$; hence if a bounded $W\in C^1$ solved
\eqref{eq:2dW}, then $F=\rho_c^2W(\rho_c,\cdot)$ would be bounded, while the
$2\pi$-periodicity of the integrand $F'$ would give
$F(\zeta+2\pi n)=F(\zeta)+nP(c)$ for every $n$, which is unbounded. This rules out every bounded $W$; the passage to a bounded smooth
weighted-divergence-free $V$ is the representation at the end of the proof of
\cref{prop:product}.
\end{proof}

\begin{remark}[classes covered by \cref{prop:periodicflat}]\label{rem:periodicflat}
Both conditions are one-dimensional tail conditions on $\varphi$, and neither
involves the order of the minima.  Written in the variable
$L=\log(1/\varepsilon)$, (i) allows $\mathcal A$ to grow like $L^{2-\eta}$
(it is deliberately weaker than boundedness of $(\rho A)_-$, which is
$\eta=2$), while (ii) allows the curvature weight to vanish like
$-c_*L^{-q}\varepsilon^{-2}$ (the case $q=0$ is
$-\mathcal W'\varepsilon^2\ge c_*>0$, its original form).  The positive
curvature term of $I_2$ then grows like $L^{2-q}$ and the negative bound for
$I_1$ like $L^{2-\eta}$; thus $q<\eta$ is exactly the condition that the
curvature term wins, and it costs nothing in the families below.

For $\varphi_1=\rho^2/(1+\rho^2)$ one has $\Psi=1/(2\varepsilon^2)$,
$\mathcal W=(1-\varepsilon)/(2\varepsilon^2)$ and
\begin{equation}
\mathcal W'(\varepsilon)=-\frac1{\varepsilon^3}+\frac1{2\varepsilon^2}
\le-\frac12\varepsilon^{-2}\qquad(0<\varepsilon\le1),
\end{equation}
while $\rho A=-4\rho^2/(1+\rho^2)\ge-4$; so (i) holds with $\eta=2$, (ii)
with $q=0$ and $c_*=1/2$, and
\cref{prop:periodicflat} recovers \cref{prop:periodic} by a different
argument, valid for every positive nonconstant $2\pi$-periodic $h\in C^2(\R)$.
The genuinely new cases are the following.
\begin{itemize}
\item Power tail $1-\varphi\sim C\rho^{-p}$: $\rho A\to-(p+2)$, so (i) holds
with $\eta=2$; $\Psi\sim\widehat C\varepsilon^{-(p+2)/p}$ and
$-\mathcal W'\varepsilon^2\sim\frac{p+2}{p}\widehat
C\varepsilon^{-2/p}\to\infty$, so (ii) holds with $q=0$.
\item Exponential tail $1-\varphi=e^{-\kappa\rho}(1+\kappa\rho)$: here
$\rho A\equiv-\kappa\rho$ is unbounded, but $\varepsilon=e^{-\kappa\rho}(1+\kappa\rho)$
gives $\kappa\rho=L+O(\log L)$ with $L=\log(1/\varepsilon)$, so (i) holds with
$\eta=1$; moreover $\Psi=e^{\kappa\rho}/\kappa^2\asymp\varepsilon^{-1}$ and
\begin{equation}
-\mathcal W'\varepsilon^2=\frac{1+\kappa\rho}{\kappa^2}
\Bigl(1+\frac{1-\varepsilon}{\kappa\rho}\Bigr)\sim\frac{\rho}{\kappa}\to\infty ,
\end{equation}
so (ii) holds with $q=0$ and any $c_*$.
\item Stretched-exponential tail $\varepsilon\sim e^{-\rho^\alpha}$:
$\rho A=(\alpha-2)-\alpha\rho^\alpha$, so $(\rho A)_-=\alpha L+O(1)$ and (i)
holds with $\eta=1$ for \emph{every} $\alpha>0$.  Also
$\Psi\sim\alpha^{-1}\varepsilon^{-1}\rho^{2-\alpha}$ and
\begin{equation}
-\mathcal W'\varepsilon^2=\frac{\rho^{2-\alpha}}{\alpha}
+\frac{(1-\varepsilon)(2-\alpha)}{\alpha^2}\rho^{2-2\alpha},
\end{equation}
which tends to $+\infty$ when $\alpha<2$, equals $1/2$ when $\alpha=2$, and
decays like $\frac1\alpha\rho^{2-\alpha}=\frac1\alpha
L^{(2-\alpha)/\alpha}$ when $\alpha>2$ (the second term being negligible
against the first).  With $q(\alpha)=(\alpha-2)_+/\alpha\in[0,1)$ this is
precisely (ii): for $\alpha\le2$ the weight does not decay at all ($q=0$),
and for $\alpha>2$ it decays only like a power of $L$.  Hence (ii) holds for
\emph{every} $\alpha>0$, and since $q(\alpha)<1=\eta$ in all cases the
comparison $q<\eta$ is satisfied.  (For $\alpha>2$ the constant furnished by
(ii) is naturally small, so the divergence asserted by
\cref{prop:periodicflat} is far beyond the range of direct numerics; the claim
is asymptotic.)
\end{itemize}
Thus every profile in the three families above satisfies (i) and (ii) for
\emph{arbitrary} $C^2$ periodic $h$, flat minima included.  What remains open
is an $\mathcal A$ growing like $L^2$ or faster in (i), or a curvature weight
vanishing faster than any power of $L$ in (ii) --- for instance
$-\mathcal W'\varepsilon^2\lesssim e^{-(\log\frac1\varepsilon)^\beta}$ with
$\beta>1$ (for $\beta\le1$ this is not faster than any power of $L$).
\end{remark}
\section{The obstruction}\label{sec:obstruction}

\begin{remark}[the bounded class: obstruction and a non-admissible counterexample]\label{rem:2dopen}
By \cref{prop:shear,prop:zeroadv,prop:product,prop:min,prop:periodic,prop:periodicgen,prop:periodicflat},
the shear and zero-advection classes, the product classes covered by
\cref{prop:product,prop:min}, the entire periodic product class for
$\varphi_1$ with arbitrary $C^2$ data, the periodic product class for every
saturating $\varphi$ satisfying \eqref{eq:Htheta} with finitely many
finite-order global minima, and the periodic product class for every
saturating $\varphi$ satisfying (i)--(ii) of \cref{prop:periodicflat} with
arbitrary $C^2$ data are rigid
without any decay hypothesis. In particular the all-degenerate example
$h=1-\cos\zeta+\cos^3\zeta/3$, for which $h''=0$ at every critical point, is
now excluded for $\varphi_1$ and for every $\varphi$ satisfying
\eqref{eq:Htheta} with $K=2$. The periodic product class for a general
$\varphi$, and the general class $V^\rho\not\equiv0$, remain open; no
admissible counterexample is known.
The obstruction is visible in the exact weighted identity, valid for every
bounded smooth $D\Subset\{\rho>0\}\times\R$,
\begin{equation}\label{eq:2denergy}
\nu\int_D\frac{|\grad G|^2}{\rho}
=\nu\oint_{\partial D}\frac{G\,\partial_nG}{\rho}
-\frac12\oint_{\partial D}\frac{(V\cdot n)G^2}{\rho}
-\int_D\frac{V^\rho G^2}{\rho^2},
\end{equation}
obtained by testing \eqref{eq:2d} against $G/\rho$; here $\Delta_\mu$ is
self-adjoint in $L^2(\rho^{-1}d\rho\,d\zeta)$, and the last term is the
sign-indefinite bulk flux. It has no gradient factor and is not controlled by
the pointwise bound $|V|\le C_0$ alone; the counterexample below shows that
boundedness and axis regularity without weighted incompressibility do not
suffice. Whether the bulk flux can be dominated in the bounded class, where
the surface terms in \eqref{eq:2denergy} are also uncontrolled, is the open
question. Equivalently, with
$h=G/\rho^2$,
\begin{equation}\label{eq:2dlift}
\Div_5\bigl(\nu\grad h-Vh\bigr)=0,\qquad
\Div_5V=\frac{2V^\rho}{\rho},
\end{equation}
so $V^\rho\not\equiv0$ makes the zero-order coefficient sign-indefinite and
the maximum principle unavailable. In the Navier--Stokes class
$V_0=2V^\rho/\rho=\partial_\zeta W$ is a $\zeta$-derivative of a bounded
function. This relation is, however, automatic for every
bounded smooth weighted-div-free field: with
$\psi=\int_0^\rho rV^\zeta(r,\zeta)\,dr$ one has
$|W|\le\|V^\zeta\|_\infty$ and $V^\rho=\rho W_\zeta/2$. It therefore carries
no sign and does not by itself close the lemma; the genuinely
Navier--Stokes-specific inputs are $(H_2)$ and the momentum coupling, neither
of which is visible in the abstract problem \eqref{eq:2d}.

The level-set form of \eqref{eq:2d} is also exact. With
$X=(G_\zeta,-G_\rho)$, \eqref{eq:2d} is the transport equation
\begin{equation}\label{eq:2dtransport}
X\cdot\grad W+\frac{2G_\zeta}{\rho}W
=-\frac{2\nu}{\rho}\Delta_\mu G .
\end{equation}
Along the characteristics $d\rho/ds=G_\zeta$, $d\zeta/ds=-G_\rho$ the
integrating factor is $\rho^2$, so on a closed level curve $\gamma$
\begin{equation}\label{eq:2dfredholm}
\oint_\gamma\rho\,\Delta_\mu G\,ds=0 .
\end{equation}
This is an identity, not a constraint: on-shell $\rho^2W=-2\psi$ is
single-valued, so the circulation in \eqref{eq:2dfredholm} vanishes
automatically. The only local content is $\Delta G(P)=0$ at a strict interior
extremum $P$. Closed-orbit solvability therefore cannot supply the missing
sign in \eqref{eq:2denergy}. The closedness of $\gamma$ is essential: a
periodic open level graph in the plane is not a closed loop, and the period
integral need not vanish (\cref{rem:drift}).

Boundedness, axis regularity and $G(0,\zeta)=0$ alone do not suffice: the
field
\begin{equation}\label{eq:2dcounter}
V^\rho=-\frac{4\nu\rho}{1+\rho^2},\qquad V^\zeta=0,\qquad
G=\frac{\rho^2}{1+\rho^2}
\end{equation}
satisfies $V\cdot\grad G=\nu\Delta_\mu G$ exactly, with $V$ bounded,
$V^\rho=O(\rho)$, $G\in C^2\cap L^\infty$ and
$G(0)=\partial_\rho G(0)=0$, yet $G\not\equiv0$. It is \emph{not}
admissible: its weighted divergence
$\partial_\rho(\rho V^\rho)+\partial_\zeta(\rho V^\zeta)
=-8\nu\rho/(1+\rho^2)^2\ne0$, and any
weighted-div-free completion would require
$V^\zeta=8\nu\zeta/(1+\rho^2)^2+O(1)$, which is unbounded in $\zeta$. The
mechanism is escape to $\rho=\infty$: the scale density of the
$\rho$-marginal is $\rho/(1+\rho^2)^2$, whose total mass is finite, so the
axis is reached only with probability $1/(1+\rho_0^2)<1$. Weighted
incompressibility is therefore essential. The surviving open case is the
sign-indefinite flux in \eqref{eq:2denergy}, of the same sign-indefinite,
gradient-free type as the averaged correlation
$C=\langle b^\rho\Gamma\rangle$ of the companion paper \cite{companion}.

Finally, the axis condition $G=O(\rho^2)$ cannot be dropped: with $V=0$,
the $\Delta_\mu$-harmonic profile
$G=\rho^2(\rho^2+\zeta^2)^{-3/2}$ is bounded at spatial infinity but blows
up at the origin along $\zeta=0$.
\end{remark}


\section*{Funding}
This research received no specific grant from any funding agency in the
public, commercial, or not-for-profit sectors.

\section*{Conflict of interest}
The author declares that there is no conflict of interest.

\begin{thebibliography}{9}
\bibitem{companion}
\emph{Companion paper}: the conditional Liouville theorem for ancient
axisymmetric Navier--Stokes flow with swirl at the critical endpoint.
% TODO: replace with the final reference once Paper A is submitted.
\bibitem{KNSS}
G.~Koch, N.~Nadirashvili, G.~Seregin, V.~\v{S}ver\'ak,
\emph{Liouville theorems for the Navier--Stokes equations and applications},
Acta Math. \textbf{203} (2009), 83--105.
\bibitem{Doob}
J.~L. Doob, \emph{Classical Potential Theory and Its Probabilistic
Counterpart}, Grundlehren der mathematischen Wissenschaften \textbf{262},
Springer, 1984.
\bibitem{Pinsky}
R.~G. Pinsky, \emph{Positive Harmonic Functions and Diffusion},
Cambridge Studies in Advanced Mathematics \textbf{45}, Cambridge University
Press, 1995.
\bibitem{SSSZ}
G.~Seregin, L.~Silvestre, V.~\v{S}ver\'ak, A.~Zlato\v{s}, \emph{On
divergence-free drifts}, J. Differential Equations \textbf{252} (2012), no.~1,
505--540; arXiv:1010.6025.
\bibitem{LRZperiodic}
Z.~Lei, X.~Ren, Q.~S. Zhang, \emph{A Liouville theorem for axi-symmetric
Navier--Stokes equations on $\R^2\times\mathbb{T}^1$}, Math. Ann.
\textbf{383} (2022), 415--431.
\bibitem{LZZ}
Z.~Lei, Q.~S. Zhang, N.~Zhao, \emph{Improved Liouville theorems for axially
symmetric Navier--Stokes equations}, Scientia Sinica Mathematica \textbf{47}
(2017), no.~10, 1183--1198; arXiv:1701.00868.
\end{thebibliography}

\end{document}
